\documentclass{article}
\usepackage{amsmath,amssymb,amsthm}
\usepackage{fancyhdr}
\usepackage{bm}
\usepackage{amsfonts,amssymb,latexsym}
\usepackage{graphicx}
\usepackage{enumitem}
\DeclareFontFamily{U}{eul}{}
\DeclareFontShape{U}{eul}{b}{n}{ <5> <6> <7> <8> gen * eurm
  <10> <10.95> <12> <14.4> <17.28> <20.74> <24.88> eurm10}{}
\DeclareSymbolFont{euler}{U}{eul}{b}{n}
\DeclareMathSymbol{\bfbeta}{\mathalpha}{euler}{'014}
\DeclareMathSymbol{\bfeps}{\mathalpha}{euler}{'017}
\DeclareFontShape{U}{cmr}{b}{n}{ <5> <6> <7> <8> gen * cmr
  <10> <10.95> <12> <14.4> <17.28> <20.74> <24.88> cmr10}{}
\DeclareSymbolFont{cmrb}{U}{cmr}{b}{n}
\DeclareMathSymbol{\omeg}{\mathalpha}{cmrb}{'012}
\setlength{\topmargin}{0in}
\setlength{\oddsidemargin}{0in}
\setlength{\textwidth}{6.5in}
\setlength{\textheight}{8in}
\newcommand{\Var}{\mathrm{Var}}
\newcommand{\op}{o_p}
\newcommand{\pconv}{\xrightarrow{p}}
\newcommand{\dconv}{\xrightarrow{d}}
\newcommand{\ind}{\perp\!\!\!\perp}
\newcommand{\E}{\mathbb{E}}
\newcommand{\tauate}{\tau}
\theoremstyle{definition}
\newtheorem{problem}{Problem}
\newtheorem{subproblem}{Part}[problem]
\input{stat.sty}
\pagestyle{fancy}
\fancyhf{}
\lhead{Gulotty}
\chead{Political Science 307}
\rhead{Assignment \#4}
\lfoot{\today}
\rfoot{Page \thepage}

\begin{document}

This assignment covers asymptotic theory, maximum likelihood and Probit estimation, hypothesis testing, and the theory of efficient estimation.  Problems 1--3 develop and apply the core asymptotic tools.  Problems 4--9 build the theory of influence functions and doubly-robust estimation, continuing the simulation from HW2.  Show your work.

\bigskip

\begin{enumerate}

%%% PROBLEM 1 %%%
\item \textbf{Convergence and the Delta Method.}

\begin{enumerate}
  \item Given $X_n \sim N(\mu, \sigma^2/n)$, prove convergence in probability to $\mu$ using Chebyshev's inequality. That is, show that for every $\epsilon > 0$,
  \[
    \Pr\!\left(|X_n - \mu| > \epsilon\right) \to 0 \quad \text{as } n \to \infty.
  \]

  \item For $g(x) = e^x$, derive the asymptotic distribution of $g(X_n)$ using the Delta Method. That is, show that
  \[
    \sqrt{n}\!\left(e^{X_n} - e^{\mu}\right) \xrightarrow{d} N\!\left(0,\; \sigma^2 e^{2\mu}\right).
  \]

\end{enumerate}

%%% PROBLEM 2 %%%
\item \textbf{Item Response Theory as a Probit Model.}

Install the \texttt{unvotes} package. This package provides data on United Nations General Assembly voting from 1946 to the present (Voeten, ``Data and Analyses of Voting in the UN General Assembly,'' \emph{Routledge Handbook of International Organization}, 2013). It contains three datasets:
\begin{itemize}
  \item \texttt{un\_votes}: country-level votes (``yes,'' ``no,'' or ``abstain'') on each roll call
  \item \texttt{un\_roll\_calls}: metadata for each roll call (date, session, description)
  \item \texttt{un\_roll\_call\_issues}: issue-area classifications (e.g., Arms control, Palestinian conflict, Human rights, etc.)
\end{itemize}

The \textbf{Rasch (1PL) IRT model} specifies that country $i$ votes ``yes'' on resolution $j$ with probability
\[
  P(Y_{ij} = 1) \;=\; \Phi(\theta_i - \beta_j),
\]
where $\theta_i$ is the country's \emph{ideal point} and $\beta_j$ is the resolution's \emph{difficulty}. This is nothing more than a probit regression with country and resolution fixed effects.

\begin{enumerate}
  \item \textbf{Data preparation.} Merge the three datasets. Filter to a single recent session (e.g., session~72) and to roll calls that have an issue-area classification. Create a binary outcome $Y_{ij} = 1$ if country $i$ voted ``yes'' on resolution $j$, and $Y_{ij} = 0$ otherwise (combining ``no'' and ``abstain''). Select approximately 20 countries: include the P5 (China, France, Russia, United Kingdom, United States) plus a geographically diverse set (e.g., Brazil, India, Japan, Nigeria, Egypt, Cuba, Israel, Australia, South Africa, Mexico, Germany, Saudi Arabia, Turkey, South Korea, New Zealand). Keep the data in long format (one row per country--resolution pair). Report the dimensions of your final dataset.

  \item \textbf{IRT as probit---model and identification.} Write down the Rasch model $P(Y_{ij}=1)=\Phi(\theta_i - \beta_j)$ and show that it is equivalent to the probit regression
  \[
    \texttt{glm(vote\_yes \char`~ country + rcid, family = binomial(link = "probit"))}.
  \]
  Specifically: what plays the role of the regressor matrix $X$? What plays the role of the coefficient vector? What scale normalization does the probit link impose (i.e., what is $\sigma$)? One country (or resolution) must be dropped as a reference category---what does this mean for the interpretation of the estimated ideal points?

  \item \textbf{Estimation and interpretation.} Fit the probit model above. Extract the country coefficients as estimated ideal points $\hat\theta_i$. Plot them (e.g., a dot plot or number line). Do the positions match known geopolitical alignments (e.g., US vs.\ Russia, EU clustering)? Then fit a \textbf{linear probability model} of the same specification. Are the LPM country fixed effects similar to the probit ideal points in \emph{ranking}?

  \item \textbf{Robust inference.} Compute model-based standard errors (from \texttt{summary(glm(...))}) and sandwich standard errors (using \texttt{vcovHC} from the \texttt{sandwich} package) for the country coefficients. Compare them. The probit model-based SEs assume conditional independence of votes given the fixed effects. Is this plausible? What kind of dependence might exist across resolutions within a session? (Connect to the sandwich/QMLE discussion from lecture.)
\end{enumerate}

\pagebreak

%%% PROBLEM 3 %%%
\item \textbf{Hypothesis Testing and Robust Inference.}

Install the \texttt{faraway} package and load the data with \texttt{data(teengamb)}.

There are 47 observations from a survey (Ide-Smith and Lea, \emph{Journal of Gambling Behavior}, vol.\ 4, 1988, pp.\ 110--118) with the following variables:
\begin{itemize}
  \item \emph{sex}, where $1 =$ female
  \item \emph{status}: Socioeconomic status score based on parents' occupational income
    \item \emph{income}: pounds per week
  \item \emph{verbal}: verbal score in words out of 12 correctly defined
  \item \emph{gamble}: expenditures on gambling (in pounds)
\end{itemize}

\begin{enumerate}
  \item Fit a model for gambling expenditures with the remaining variables as predictors. Report the results in a table suitable for publication.

  \item Interact gender with verbal score. Compute an $F$-test on the joint hypothesis that the coefficients on gender and gender $\times$ verbal score are both zero.

  \item Compute a Wald test of the same hypothesis using HC2 standard errors. Show that the $F$-test and Wald test are equivalent under normality and homoskedasticity; discuss when they diverge asymptotically.

  \item For the single-coefficient case, show algebraically that $t^2 = F$. That is, show that the square of the $t$-statistic for testing $H_0\colon \beta_j = 0$ equals the $F$-statistic from a regression with and without $x_j$.
\end{enumerate}

%%% Problem 4: Simulation Exercise Part II %%%
\item \textbf{Simulation Exercise Part II: Different Estimators, Same First-Order Behavior}

This problem continues the Monte Carlo experiment from HW2 Problem 4.

\textbf{Monte Carlo.} Using the same DGP as in HW2 Problem 4, compute the following additional ATE estimators for $n\in\{200,1000\}$ with at least $B=2000$ replications:

\textbf{(A) Difference in means:} $\widehat{\tau}_{\mathrm{DM}} = \bar Y_1 - \bar Y_0$.

\textbf{(B) IPW with known propensity:} Because treatment is randomized, $e(X_i)=p$ is known. Define
\[
\widehat{\tau}_{\mathrm{IPW}}
=
\frac{1}{n}\sum_{i=1}^n
\left(\frac{D_i Y_i}{p} - \frac{(1-D_i)Y_i}{1-p}\right).
\]

\textbf{(C) Regression adjustment (imputation):} Estimate $\widehat{\mu}_d(x)$ by a linear regression of $Y$ on $X$ within group $D=d$, then form
\[
\widehat{\tau}_{\mathrm{RA}}
=
\frac{1}{n}\sum_{i=1}^n \left(\widehat{\mu}_1(X_i)-\widehat{\mu}_0(X_i)\right).
\]

\textbf{(D) Lin estimator (interacted):} As in HW2 Problem 4.

\begin{enumerate}
\item \textbf{Variance ranking.} For each $n$, rank the estimators by Monte Carlo variance. Are (C) and (D) close?

\item \textbf{Estimator differences.} For each replication, compute $\widehat{\tau}_{\mathrm{Lin}} - \widehat{\tau}_{\mathrm{RA}}$ and $\widehat{\tau}_{\mathrm{Lin}} - \widehat{\tau}_{\mathrm{DM}}$. Report the Monte Carlo mean and standard deviation of these differences for each $n$. Do these differences appear to shrink as $n$ increases?
\end{enumerate}

\textbf{Theory.}

Define $\mu_d(x)=\E[Y(d)\mid X=x]$ and the residuals $\varepsilon_{i,d} = Y_i(d) - \mu_d(X_i)$ with $\E[\varepsilon_{i,d}\mid X_i]=0$. Consider the infeasible estimator
\[
\widehat{\tau}_\star
=
\frac{1}{n}\sum_{i=1}^n
\left[
\mu_1(X_i)-\mu_0(X_i)
+
\frac{D_i}{p}\big(Y_i-\mu_1(X_i)\big)
-
\frac{1-D_i}{1-p}\big(Y_i-\mu_0(X_i)\big)
\right].
\]

\begin{enumerate}
\setcounter{enumi}{2}
\item \textbf{Unbiasedness.} Show that $\E[\widehat{\tau}_\star]=\tau$.

\item \textbf{Centering.} Show that
\[
\widehat{\tau}_\star-\tau
=
\frac{1}{n}\sum_{i=1}^n \psi(Z_i),
\quad\text{where}\quad
\E[\psi(Z_i)]=0.
\]
Write $\psi(Z)$ explicitly.\footnote{Here $Z_i$ is simply shorthand for the full observed data on unit $i$,
$Z_i=(Y_i,D_i,X_i)$.
We introduce this notation so that expressions like
$\frac{1}{n}\sum_{i=1}^n \psi(Z_i)$
can be recognized as averages of i.i.d.\ random variables, to which the law of large numbers and central limit theorem apply.}

\item \textbf{CLT for the infeasible estimator.}
Using your result from part~(d), show that the infeasible estimator satisfies a CLT:
\[
  \sqrt{n}\bigl(\widehat{\tau}_\star - \tau\bigr)
  \;\dconv\;
  \mathcal{N}\!\bigl(0,\;\Var(\psi(Z))\bigr).
\]
State clearly which classical theorem you are invoking and why its
assumptions hold.

\end{enumerate}

\end{enumerate}

\bigskip

The remaining problems show that replacing the true $\mu_d$ with an estimate $\widehat{\mu}_d$ does not change the limiting distribution---a property called \emph{asymptotic equivalence}.  The \textbf{feasible} AIPW estimator $\widehat{\tau}(\widehat{\mu})$ has the same form as $\widehat{\tau}_\star$ but with each $\mu_d$ replaced by an estimator $\widehat{\mu}_d$.

\subsection*{Reference: stochastic order notation}

In asymptotic statistics we need a way to say that a random sequence
``vanishes'' or ``stays bounded'' as $n\to\infty$, even though the
quantities are random.  The two workhorses are $\op$ (little-$o$ in
probability) and $O_p$ (big-$O$ in probability).

\medskip
\textbf{Definition.}
Let $\{W_n\}$ be a sequence of random variables and $\{a_n\}$ a sequence
of positive constants.
\begin{itemize}[nosep]
  \item $W_n = \op(a_n)$ means
    $W_n/a_n \pconv 0$, i.e.\ for every $\epsilon>0$,
    \[
      P\!\bigl(|W_n/a_n|>\epsilon\bigr) \;\to\; 0
      \quad\text{as } n\to\infty.
    \]
    Read this as ``$W_n$ is negligible relative to $a_n$.''

  \item $W_n = O_p(a_n)$ means $W_n/a_n$ is
    \emph{bounded in probability}: for every $\epsilon>0$ there exists
    $M<\infty$ such that
    $P\!\bigl(|W_n/a_n|>M\bigr)<\epsilon$ for all large~$n$.
\end{itemize}

\medskip
\textbf{Key examples.}
\begin{enumerate}[nosep,label=(\roman*)]
  \item If $W_n\pconv 0$, then $W_n = \op(1)$.
  \item By the law of large numbers,
    $\bar{X}_n - \mu = \op(1)$.
    By the CLT, $\bar{X}_n - \mu = O_p(n^{-1/2})$.
  \item Saying a remainder is $\op(n^{-1/2})$ means that after
    multiplying by $\sqrt{n}$, it vanishes---so it cannot affect the
    limiting distribution of anything that is $O_p(n^{-1/2})$.
\end{enumerate}

\medskip
\textbf{Useful rules} (you may use these without proof).
\begin{enumerate}[nosep,label=(\roman*)]
  \item $\op(a_n)+\op(a_n)=\op(a_n)$.
  \item $O_p(a_n)\cdot\op(b_n) = \op(a_n b_n)$.
\item Slutsky's Lemma:
If $X_n \xrightarrow{d} X$ and $R_n \xrightarrow{p} 0$,
then $X_n + R_n \xrightarrow{d} X$.
Equivalently, if $R_n = o_p(1)$, then the same result holds.

\end{enumerate}

\medskip
\textbf{A common proof strategy.}
To show $W_n = \op(a_n)$, it often suffices to show
$\E[W_n/a_n]=0$ and $\Var(W_n/a_n)\to 0$, then apply Chebyshev:
\[
  P\!\bigl(|W_n/a_n|>\epsilon\bigr)
  \;\leq\;
  \frac{\Var(W_n/a_n)}{\epsilon^2}
  \;\to\; 0.
\]
This is exactly the strategy you will use in Problem~\ref{prob:remainder}.

\bigskip
\hrule
\bigskip

\begin{enumerate}\setcounter{enumi}{4}

%%% PROBLEM 5 %%%
\item \textbf{Algebraic decomposition of the remainder.}\label{prob:algebra}
  Define the estimation error $\Delta_d(x) = \widehat{\mu}_d(x) - \mu_d(x)$ for $d \in \{0,1\}$.

  \begin{enumerate}[label=(\alph*)]
    \item Write down the feasible estimator $\widehat{\tau}(\widehat{\mu})$
      explicitly (it has the same form as $\widehat{\tau}_\star$ with
      $\widehat{\mu}_d$ in place of $\mu_d$).

    \item By subtracting $\widehat{\tau}_\star$ from
      $\widehat{\tau}(\widehat{\mu})$, show that the $Y_i$ terms cancel
      within each treatment arm, and obtain
      \[
        \widehat{\tau}(\widehat{\mu}) - \widehat{\tau}_\star
        \;=\;
        \frac{1}{n}\sum_{i=1}^n
        \left[
          \Bigl(1-\frac{D_i}{p}\Bigr)\Delta_1(X_i)
          \;-\;
          \Bigl(1-\frac{1-D_i}{1-p}\Bigr)\Delta_0(X_i)
        \right].
      \]
  \end{enumerate}

\bigskip

%%% PROBLEM 6 %%%
\item \textbf{Neyman orthogonality.}\label{prob:orthogonality}
  This problem establishes the key property that makes AIPW ``forgiving''
  of estimation error in $\widehat\mu_d$.

  \begin{enumerate}[label=(\alph*)]
    \item Using the fact that $D_i \ind X_i$ with $P(D_i=1)=p$, show that
      \[
        \E\!\left[1-\frac{D_i}{p}\;\Big|\; X_i\right] = 0.
      \]

    \item Explain in one sentence why this means that, \emph{for any fixed
      function} $\Delta_1(\cdot)$,
      \[
        \E\!\left[\Bigl(1-\frac{D_i}{p}\Bigr)\Delta_1(X_i)\;\Big|\;X_i\right]
        = 0.
      \]

    \item Compute
      $\E\!\bigl[(1-D_i/p)^2 \mid X_i\bigr]$.
      \emph{Hint: $\Var(D_i/p\mid X_i)=(?)$.}
  \end{enumerate}

\bigskip

%%% PROBLEM 7 %%%
\item \textbf{Controlling the remainder.}\label{prob:remainder}
  We now show that the remainder from Problem~\ref{prob:algebra} is
  $\op(n^{-1/2})$ under mild conditions.
  It suffices to treat the $d=1$ term; the $d=0$ term is analogous.

  \medskip\noindent
  \textbf{Assumption (Sample splitting).}
  The sample is split into two independent halves.
  $\widehat{\mu}_1$ is estimated on the first half, and the sum in
  Problem~\ref{prob:algebra} runs over the second half of size $n$.
  (In practice one uses cross-fitting; here we keep things simple.)

  \medskip\noindent
  \textbf{Assumption ($L^2$-consistency).}
  $\|\widehat{\mu}_1-\mu_1\|_{L^2}^2
  \;=\; \E\bigl[(\widehat{\mu}_1(X)-\mu_1(X))^2\bigr]
  \;\pconv\; 0.$

  \medskip
  Define
  $R_n = \frac{1}{n}\sum_{i=1}^{n}
    \bigl(1-D_i/p\bigr)\,\Delta_1(X_i)$,
  where the sum is over the estimation (second) half of the sample.

  \begin{enumerate}[label=(\alph*)]
    \item Conditional on the training fold (hence on $\widehat{\mu}_1$),
      the summands in $R_n$ are i.i.d.  Using your results from
      Problem~\ref{prob:orthogonality}, show that $R_n$ has conditional
      mean zero and conditional variance
      \[
        \Var(R_n \mid \widehat{\mu}_1)
        \;=\;
        \frac{1-p}{n\,p}\;
        \|\widehat{\mu}_1-\mu_1\|_{L^2}^2.
      \]

    \item Apply Chebyshev's inequality \emph{conditionally}: for any
      $\epsilon>0$,
      \[
        P\!\left(\bigl|\sqrt{n}\,R_n\bigr|>\epsilon
          \;\Big|\;\widehat{\mu}_1\right)
        \;\leq\; \frac{n\;\Var(R_n\mid\widehat{\mu}_1)}{\epsilon^2}.
      \]
      Substitute part~(a) and simplify.  Note how the factor of $n$
      cancels.

    \item Using the $L^2$-consistency assumption, conclude that the
      conditional bound vanishes in probability.  Then take (outer)
      expectations to deduce $\sqrt{n}\,R_n = \op(1)$, i.e.,
      $R_n = \op(n^{-1/2})$.
  \end{enumerate}

\bigskip

%%% PROBLEM 8 %%%
\item \textbf{Putting it all together.}\label{prob:clt}
  Combine your results from the previous problems to conclude that, under
  sample splitting and $L^2$-consistency of both $\widehat{\mu}_1$ and
  $\widehat{\mu}_0$,
  \[
    \sqrt{n}\bigl(\widehat{\tau}(\widehat{\mu})-\tau\bigr)
    \;\dconv\;
    \mathcal{N}\!\bigl(0,\;\Var(\psi(Z))\bigr).
  \]
  State explicitly which results you rely upon.

\bigskip

%%% PROBLEM 9 %%%
\item \textbf{Reflection (no computation required).}\label{prob:reflect}
  \begin{enumerate}[label=(\alph*)]
    \item The $L^2$-consistency assumption imposes \emph{no rate} on how
      fast $\widehat{\mu}_d$ converges.  Why is this remarkable compared
      to a plug-in estimator that simply uses
      $n^{-1}\sum_i[\widehat{\mu}_1(X_i)-\widehat{\mu}_0(X_i)]$?

    \item What role does the Neyman-orthogonality property
      (Problem~\ref{prob:orthogonality}(a)) play in making this possible?
      In particular, what would go wrong if the conditional mean of the
      remainder were nonzero?
  \end{enumerate}
\end{enumerate}
\end{document}
